\Needspace{9\baselineskip}\begin{center}\textit{Theory of the roots of the equation \(x^n-1=0\) (where \(n\) is assumed prime).} \textit{Omitting the root \(1\), the remaining roots \((\Omega)\) are contained in the equation \(X=x^{n-1}+x^{n-2}+\mathrm{etc.}+x+1=0\).} \textit{The function \(X\) cannot be resolved into factors of lower degree whose coefficients are all rational.}\end{center}


\Needspace{6\baselineskip}\Article{341}
\textsc{Theorem.} \textit{If the function \(X\) is divisible by the function of lower degree} \[P=x^\lambda+Ax^{\lambda-1}+Bx^{\lambda-2}+\ldots+Kx+L\] \textit{the coefficients \(A,\allowbreak B\ldots L\) cannot all be integers.}

\textit{Proof.} Let \(X=PQ\), and let \(\mathfrak P\) be the collection of roots of the equation \(P=0\), and \(\mathfrak Q\) the collection of roots of the equation \(Q=0\), so that \(\Omega\) consists of \(\mathfrak P\) and \(\mathfrak Q\) taken together. Further, let \(\mathfrak R\) be the collection of roots reciprocal to those in \(\mathfrak P\), and \(\mathfrak S\) the collection of roots reciprocal to those in \(\mathfrak Q\). Let the roots contained in \(\mathfrak R\) be the roots of the equation \(R=0\) (which is easily seen to be \(x^\lambda+\frac KLx^{\lambda-1}+\mathrm{etc.}+\frac ALx+\frac1L=0\)), and those contained in \(\mathfrak S\) the roots of the equation \(S=0\). Clearly the roots \(\mathfrak R\) and \(\mathfrak S\) taken together also form the collection \(\Omega\), and \(RS=X\). We now distinguish four cases.
\SourceLeaf{423}

I. When \(\mathfrak P\) coincides with \(\mathfrak R\), and hence \(P=R\). In this case the roots in \(\mathfrak P\) clearly come in reciprocal pairs, so \(P\) is a product of \(\frac12\lambda\) quadratic factors of the form \(xx-2x\cos\omega+1\); since such a factor is \(=(x-\cos\omega)^2+\sin\omega^2\), it is easily seen that \(P\) necessarily takes a positive real value for every real value of \(x\). Let the equations whose roots are respectively the squares, cubes, fourth powers \(\ldots\) and powers of order \(n-1\) of the roots in \(\mathfrak P\) be \(P^{\prime}=0\), \(P^{\prime\prime}=0,\allowbreak P^{\prime\prime\prime}=0,\allowbreak \ldots P^\nu=0\), and let the values of the functions \(P,\allowbreak P^{\prime},\allowbreak P^{\prime\prime}\ldots P^\nu\) obtained by setting \(x=1\) be respectively \(p,\allowbreak p^{\prime},\allowbreak p^{\prime\prime}\ldots p^\nu\). Then, by what was said above, \(p\) will be a positive quantity, and for exactly the same reason \(p^{\prime},\allowbreak p^{\prime\prime}\), etc., will also be positive. Now \(p\) is the value of the function \((1-t)(1-u)(1-v)\), etc., obtained by substituting for \(t,\allowbreak u,\allowbreak v\), etc., the roots in \(\mathfrak P\); \(p^{\prime}\) is the value of the same function when we substitute for \(t,\allowbreak u,\allowbreak v\), etc., the squares of those roots, etc.; moreover, the value for \(t=1,\allowbreak u=1,\allowbreak v=1\), etc., is clearly \(=0\). Thus the sum \(p+p^{\prime}+p^{\prime\prime}\ldots+p^\nu\) will be an integer divisible by \(n\). It is also easily seen that the product \(PP^{\prime} P^{\prime\prime}\ldots\) is \(=X^\lambda\), and hence \(pp^{\prime} p^{\prime\prime}\ldots=n^\lambda\). Now, if all the coefficients in \(P\) were rational, all those in \(P^{\prime},\allowbreak P^{\prime\prime}\), etc., would also be rational by Article 338; by Article 42, however, all these coefficients would necessarily be integers. Thus \(p,\allowbreak p^{\prime},\allowbreak p^{\prime\prime}\), etc., would all be integers; since their product is \(n^\lambda\), and their number satisfies \(n-1>\lambda\), some of them (at least \(n-1-\lambda\)) must necessarily be \(=1\), and the rest equal to \(n\) or to a power of \(n\). Thus, if \(g\) of them are \(=1\), the sum \(p+p^{\prime}+\mathrm{etc.}\) will clearly be \(\equiv g\ (\mathrm{mod.}\ n)\), and hence certainly not divisible by \(n\). The assumption therefore cannot hold.

II. When \(\mathfrak P\) and \(\mathfrak R\) do not coincide but contain some common roots, let \(\mathfrak T\) be their collection, and \(T=0\) the equation whose roots they are. Then \(T\) will be the greatest common divisor of the functions \(P,\allowbreak R\) (as is known from the theory of equations). The roots in \(\mathfrak T\) clearly come in reciprocal pairs; hence, by what has been proved above, the coefficients in \(T\) cannot all be rational. But they certainly would be if all those in \(P\), and hence also all those in \(R\), were rational, as follows directly from the nature of the operation of finding the greatest common divisor. The assumption is therefore absurd.

III. When \(\mathfrak Q\) and \(\mathfrak S\) either coincide or at least contain common roots, in exactly the same way the coefficients in \(Q\) cannot all be rational; but they would be rational if all those in \(P\) were rational. This is therefore impossible.

IV. But if neither the pair \(\mathfrak P\) and \(\mathfrak R\) nor the pair \(\mathfrak Q\) and \(\mathfrak S\) has any common\SourceLeaf{424}
root, all the roots \(\mathfrak P\) must be found in \(\mathfrak S\), and all those in \(\mathfrak Q\) in \(\mathfrak R\), whence \(P=S\) and \(Q=R\). Therefore \(X=PQ\) will be the product of \(P\) and \(R\), that is, \[\begin{gathered}
\text{of}\quad x^\lambda+Ax^{\lambda-1}\ldots+Kx+L\\
\text{and}\quad x^\lambda+\frac KLx^{\lambda-1}\ldots+\frac ALx+\frac1L
\end{gathered}\] whence, setting \(x=1\), we obtain \[nL=(1+A\ldots+K+L)^2\] Now, if all the coefficients in \(P\) were rational, and hence also integers by Article 42, \(L\), which must divide the last coefficient in \(X\), namely unity, would necessarily be \(=\pm1\), whence \(\pm n\) would be a square. Since this contradicts the hypothesis, the assumption cannot hold.

This theorem therefore shows that, however \(X\) is resolved into factors, at least some of their coefficients will be irrational, and hence cannot be determined except by an equation of higher degree.

\Needspace{9\baselineskip}\begin{center}\textit{The aim of the following investigations is explained.}\end{center}


\Needspace{6\baselineskip}\Article{342}
The aim of the following investigations, which it will be useful to explain briefly, is to resolve \(X\) \textsc{step by step} into more and more factors, determining their coefficients by equations of the lowest possible degree, until we reach linear factors, or the roots \(\Omega\) themselves. Namely, we will show that, if the number \(n-1\) is resolved in any way into integer factors \(\alpha,\allowbreak \beta,\allowbreak \gamma\), etc. (each of which may be taken prime), \(X\) can be resolved into \(\alpha\) factors of degree \(\frac{n-1}{\alpha}\), whose coefficients are determined by an equation of degree \(\alpha\); each of these factors in turn into \(\beta\) others of degree \(\frac{n-1}{\alpha\beta}\), with the help of an equation of degree \(\beta\), etc. Thus, if \(\nu\) denotes the number of factors \(\alpha,\allowbreak \beta,\allowbreak \gamma\), etc., finding the roots \(\Omega\) reduces to solving \(\nu\) equations of degrees \(\alpha,\allowbreak \beta,\allowbreak \gamma\), etc. \textit{For example}, for \(n=17\), where \(n-1=2.2.2.2\), four quadratic equations must be solved; for \(n=73\), three quadratic and two cubic equations.

In what follows we will very often have to consider powers of the root \(r\) whose exponents are themselves powers, and such expressions are troublesome to set in type. To make printing easier, we will henceforth\SourceLeaf{425}
use the following abbreviation. For \(r,\allowbreak rr,\allowbreak r^3\), etc., we will write \([1],\allowbreak [2],\allowbreak [3]\), etc., and generally for \(r^\lambda\), where \(\lambda\) denotes any integer, \([\lambda]\). Such expressions are therefore not yet fully determined, but become so as soon as, for \(r\) or \([1]\), a specific root in \(\Omega\) is chosen. Thus, in general \([\lambda],\allowbreak [\mu]\) will be equal or unequal according as \(\lambda,\allowbreak \mu\) are congruent or incongruent modulo \(n\); furthermore \([0]=1;\allowbreak \ [\lambda].[\mu]=[\lambda+\mu];\allowbreak \ [\lambda]^\nu=[\lambda\nu]\); the sum \([0]+[\lambda]+[2\lambda]\ldots+[(n-1)\lambda]\) is either \(0\) or \(n\), according as \(\lambda\) is not divisible or is divisible by \(n\).

\Needspace{9\baselineskip}\begin{center}\textit{All the roots \(\Omega\) are distributed into certain classes (periods).}\end{center}


\Needspace{6\baselineskip}\Article{343}
If, for the modulus \(n\), \(g\) is a number of the kind we called a primitive root in Section III, the \(n-1\) numbers \(1,\allowbreak g,\allowbreak gg\ldots g^{n-2}\) will be congruent to \(1,\allowbreak 2,\allowbreak 3\ldots n-1\) modulo \(n\), although in a different order: each number of one sequence has a congruent number in the other. It follows at once that the roots \([1],\allowbreak [g],\allowbreak [gg]\ldots[g^{n-2}]\) coincide with \(\Omega\); in exactly the same way, more generally, \[\relax[\lambda],[\lambda g],[\lambda gg]\ldots[\lambda g^{n-2}]\quad\text{with}\quad\Omega\] coincide, where \(\lambda\) denotes any integer not divisible by \(n\). Further, since \(g^{n-1}\equiv1\ (\mathrm{mod.}\ n)\), it is readily seen that the two roots \([\lambda g^\mu],\allowbreak [\lambda g^\nu]\) are identical or different according as \(\mu,\allowbreak \nu\) are congruent or incongruent modulo \(n-1\). Thus, if \(G\) is another primitive root, the roots \([1],\allowbreak [g]\ldots[g^{n-2}]\) will also agree with \([1],\allowbreak [G]\ldots[G^{n-2}]\) if order is disregarded. But it is also easily proved that, if \(e\) is a divisor of \(n-1\), and we set \(n-1=ef,\allowbreak \ g^e=h\), \(G^e=H\), the \(f\) numbers \(1,\allowbreak h,\allowbreak hh\ldots h^{f-1}\) will be congruent to \(1,\allowbreak H,\allowbreak H^2\ldots H^{f-1}\) modulo \(n\) (without regard to order). For suppose \(G\equiv g^\omega\ (\mathrm{mod.}\ n)\), and let \(\mu\) be any positive number \(<f\), and \(\nu\) the least residue of \(\mu\omega\ (\mathrm{mod.}\ f)\). Then \(\nu e\equiv\mu\omega e\ (\mathrm{mod.}\ n-1)\), hence \(g^{\nu e}\equiv g^{\mu\omega e}\equiv G^{\mu e}\ (\mathrm{mod.}\ n)\), or \(H^\mu\equiv h^\nu\); that is, each number of the latter sequence \(1,\allowbreak H,\allowbreak H^2\), etc., has a congruent number in the sequence \(1,\allowbreak h,\allowbreak hh\ldots\), and likewise conversely. — It is therefore clear that the \(f\) roots \([1],\allowbreak [h],\allowbreak [hh]\ldots[h^{f-1}]\) are identical with \([1],\allowbreak [H],\allowbreak [H^2]\ldots[H^{f-1}]\); more generally, it is readily seen in the same way that \[\begin{gathered}
\relax[\lambda],[\lambda h],[\lambda hh]\ldots[\lambda h^{f-1}]\\
\text{with}\quad[\lambda],[\lambda H],[\lambda H^2]\ldots[\lambda H^{f-1}]
\end{gathered}\] agree. We will denote the \textit{sum} of such \(f\) roots \([\lambda]+[\lambda h]+\mathrm{etc.}+[\lambda h^{f-1}]\), which, since it does not\SourceLeaf{426}
change when another primitive root is taken for \(g\), must be regarded as independent of \(g\), by \((f;\allowbreak \lambda)\); the \textit{collection} of the same roots we will call the \textit{period} \((f,\allowbreak \lambda)\), without regard to the order of the roots\HistoricalNote{*)}{\normalfont In what follows, we may also call the sum the numerical value of the period, or simply the period where there is no risk of ambiguity.}. — In writing such a period, it will be useful to reduce each of its roots to the simplest expression, namely to substitute for the numbers \(\lambda,\allowbreak \lambda h,\allowbreak \lambda hh\), etc., the least residues modulo \(n\); the terms of the period may also be arranged by their size, if desired.

\textit{For example}, for \(n=19\), where \(2\) is a primitive root, the period \((6,\allowbreak 1)\) consists of the roots \([1],\allowbreak [8],\allowbreak [64],\allowbreak [512],\allowbreak [4096],\allowbreak [32768]\), or \([1],\allowbreak [7],\allowbreak [8],\allowbreak [11],\allowbreak [12],\allowbreak [18]\). Similarly the period \((6,\allowbreak 2)\) consists of \([2],\allowbreak [3],\allowbreak [5],\allowbreak [14],\allowbreak [16],\allowbreak [17]\). The period \((6,\allowbreak 3)\) is found to be identical with the preceding one. The period \((6,\allowbreak 4)\) contains \([4],\allowbreak [6],\allowbreak [9],\allowbreak [10],\allowbreak [13],\allowbreak [15]\).

\Needspace{9\baselineskip}\begin{center}\textit{Various theorems on the periods of the roots \(\Omega\).}\end{center}


\Needspace{6\baselineskip}\Article{344}
The following observations about these periods immediately present themselves:

I. Since \(\lambda h^f\equiv\lambda,\allowbreak \ \lambda h^{f+1}\equiv\lambda h\), etc., \((\mathrm{mod.}\ n)\), clearly the same roots that constitute \((f,\allowbreak \lambda)\) also constitute \((f,\allowbreak \lambda h),\allowbreak (f,\allowbreak \lambda hh)\), etc.; thus, in general, if \([\lambda']\) denotes any root in \((f,\allowbreak \lambda)\), this period will be completely identical with \((f,\allowbreak \lambda')\). Hence, if two periods consisting of equally many roots (which we will call \textit{similar}) have any common root, they are clearly identical. Thus it is impossible for two roots to occur together in one period, while only one of them occurs in another similar period. It is also clear that, if two roots \([\lambda],\allowbreak [\lambda']\) belong to the same period of \(f\) terms, the value of the expression \(\frac{\lambda'}{\lambda}\ (\mathrm{mod.}\ n)\) is congruent to some power of \(h\), or we may put \(\lambda'\equiv\lambda g^{\nu e}\ (\mathrm{mod.}\ n)\).

II. If \(f=n-1,\allowbreak \ e=1\), the period \((f,\allowbreak 1)\) clearly coincides with \(\Omega\); in the other cases \(\Omega\) will be composed of \(e\) periods \((f,\allowbreak 1),\allowbreak (f,\allowbreak g),\allowbreak (f,\allowbreak gg)\ldots(f,\allowbreak g^{e-1})\). These periods are therefore all distinct, and clearly any other similar period \((f,\allowbreak \lambda)\) coincides with one of them, provided \([\lambda]\) belongs to \(\Omega\), that is, if \(\lambda\) is not divisible by \(n\). But the period \((f,\allowbreak 0)\), or \((f,\allowbreak kn)\), clearly consists of \(f\) units. It is just as easily seen\SourceLeaf{427}
that, if \(\lambda\) is any number not divisible by \(n\), the collection of \(e\) periods \((f,\allowbreak \lambda),\allowbreak (f,\allowbreak \lambda g),\allowbreak (f,\allowbreak \lambda gg)\ldots(f,\allowbreak \lambda g^{e-1})\) also coincides with \(\Omega\). — Thus, \textit{for example}, for \(n=19,\allowbreak \ f=6\), \(\Omega\) consists of the three periods \((6,\allowbreak 1),\allowbreak (6,\allowbreak 2),\allowbreak (6,\allowbreak 4)\), to one of which any other similar period, except \((6,\allowbreak 0)\), reduces.

III. If \(n-1\) is the product of three positive numbers \(a,\allowbreak b,\allowbreak c\), clearly any period of \(bc\) terms is composed of \(b\) periods of \(c\) terms: namely, \((bc,\allowbreak \lambda)\) of \((c,\allowbreak \lambda),\allowbreak (c,\allowbreak \lambda g^a),\allowbreak (c,\allowbreak \lambda g^{2a})\ldots(c,\allowbreak \lambda g^{ab-a})\), and these will therefore be said to be contained in the former. Thus, for \(n=19\), the period \((6,\allowbreak 1)\) consists of the three periods \((2,\allowbreak 1),\allowbreak (2,\allowbreak 8),\allowbreak (2,\allowbreak 7)\); the first contains the roots \(r,\allowbreak r^{18}\), the second \(r^8,\allowbreak r^{11}\), and the third \(r^7,\allowbreak r^{12}\).

